% ============================================================
% Chapter 4: Baseline Activity-Oriented Relax-and-Fix
% Thesis: Solving the FSRCPSP Using Hybrid ML Approaches
% Author: Souvik Chakraborty
% ============================================================
\chapter{Baseline: Activity-Oriented Relax-and-Fix}
\label{chap:baseline}
Before introducing the group-based method in
Chapter~\ref{chap:grouprf}, a simpler baseline is needed. This
chapter presents an activity-oriented Relax-and-Fix~(R\&F)
heuristic that adapts the unit-oriented decomposition of
\cite{franz2019fixoptimize} to the SAA-FSRCPSP, fixing one
activity per iteration in topological order. The group-based
extension builds directly on this baseline.
\input{figures/baseline_pipeline}
% ------------------------------------------------------------
\section{Overview}
\label{sec:baseline_overview}
% ------------------------------------------------------------
Relax-and-Fix is a decomposition heuristic for mixed-integer
programmes \parencite{pochet2006,Helber2010}. The idea is to
partition the binary decision variables into subsets and fix them
iteratively: in each iteration one subset is optimised while the
rest are either locked at previously found values or relaxed to
continuous variables. Because only a small number of binaries are
active at any time, each subproblem is much smaller than the full
MIP.
Within scheduling, R\&F has been used in several forms.
\cite{etminaniesfahani2024relaxsolve} proposed a time-oriented
R\&F with rolling time windows for the multi-mode RCPSP, and
\cite{bigler2024multisite} introduced a unit-oriented R\&F for
the multi-site RCPSP. The baseline here follows the unit-oriented
idea of \cite{franz2019fixoptimize} and applies it activity by
activity to the SAA-FSRCPSP.
The procedure has two phases. First, a feasible greedy solution is
constructed via the Serial Schedule Generation Scheme~(SSGS).
Second, each activity's start time is refined through a sequence of
MIP subproblems. Figure~\ref{fig:baseline_pipeline} shows the
complete pipeline.
% ------------------------------------------------------------
\section{Initial Solution via SSGS}
\label{sec:ssgs}
% ------------------------------------------------------------
The R\&F phase needs a feasible starting schedule both to
warm-start the MIP solver and to set up the initial time windows.
This starting solution comes from a variant of the Serial Schedule
Generation Scheme~(SSGS) \parencite{Kolisch1999}, adapted here to
the flexible resource profiles of the FSRCPSP. In the standard
RCPSP, durations and resource demands are fixed, so the SSGS only
has to decide \textit{when} to schedule each activity. In the
FSRCPSP the duration~$d_i$ depends on how the workload~$w_{ik}$ is
spread over time within the bounds $[\ell_{ik}, u_{ik}]$, so the
SSGS also has to decide \textit{for how long} each activity runs.
The adapted SSGS has three parts: worst-case scenario selection
(choosing which scenario workloads to use), resource profile
construction (computing the shortest feasible duration and the
per-period allocations), and the scheduling loop (placing
activities in earliest-start order while respecting capacity).
Algorithm~\ref{alg:ssgs} gives the full procedure.
\subsection{Worst-Case Scenario Selection}
\label{subsec:worst_case}
The SAA-FSRCPSP involves multiple scenarios~$\pi \in \Pi$, each
with its own workloads~$w_{ik}^{\pi}$. The SSGS needs to pick one
workload vector per activity, and since the schedule should hold up
under all scenarios, a conservative choice makes sense.
The idea is simple: for each activity~$i$, pick the scenario that
would force the longest duration. Under scenario~$\pi$ and
resource~$k$, the minimum number of periods needed (even at full
capacity) is $\lceil w_{ik}^{\pi} / u_{ik} \rceil$. The
worst-case scenario is the one where this implied duration is
largest across all resources:
\begin{equation}
    \hat{\pi}_i = \arg\max_{\pi \in \Pi}\;
    \max_{k \in K}\;
    \left\lceil \frac{w_{ik}^{\pi}}{u_{ik}} \right\rceil
    \label{eq:worst_case}
\end{equation}
The workload vector used by the SSGS for activity~$i$ is then
$\hat{w}_{ik} = w_{ik}^{\hat{\pi}_i}$ for all~$k \in K$. By
always planning for the hardest scenario per activity, the
resulting schedule tends to be conservative but robust enough to
serve as a warm-start.
\subsection{Resource Profile Construction}
\label{subsec:profile_construction}
Once the worst-case workloads~$\hat{w}_{ik}$ are fixed, the SSGS
needs a feasible duration~$D_i$ and a per-period allocation
$r_{ik}^{(t)} \in [\ell_{ik}, u_{ik}]$ for each resource~$k$ and
period~$t \in \{1, \ldots, D_i\}$, with the total adding up to the
workload: $\sum_{t=1}^{D_i} r_{ik}^{(t)} = \hat{w}_{ik}$.
If resource~$k$ were the only one, the shortest feasible duration
would be $D_k^{\min} = \lceil \hat{w}_{ik} / u_{ik} \rceil$ ---
that is, every period at full capacity. But since all resources
share one common duration, the binding one wins. On the other end,
if a lower bound $\ell_{ik} > 0$ applies, the duration cannot
exceed $D_k^{\max} = \lfloor \hat{w}_{ik} / \ell_{ik} \rfloor$
(otherwise the per-period allocation would drop below the lower
bound). So the common duration is:
\begin{equation}
    D_i = \max_{k \in K}\; D_k^{\min}
    = \max_{k \in K}\;
    \left\lceil \frac{\hat{w}_{ik}}{u_{ik}} \right\rceil
    \label{eq:duration}
\end{equation}
subject to the feasibility condition
$D_i \leq \min_{k : \ell_{ik} > 0} D_k^{\max}$.
With $D_i$ fixed, the per-period allocations are built using a
\textit{front-loaded} strategy. Every period starts at the lower
bound~$\ell_{ik}$, and whatever workload remains
($\hat{w}_{ik} - D_i \cdot \ell_{ik}$) is packed in from left to
right, filling each period up to~$u_{ik}$ before moving on.
Formally:
\begin{equation}
    r_{ik}^{(t)} = \ell_{ik}
    + \min\!\left(u_{ik} - \ell_{ik},\;
    \hat{w}_{ik} - D_i \cdot \ell_{ik}
    - \sum_{\tau=1}^{t-1}
    \bigl(r_{ik}^{(\tau)} - \ell_{ik}\bigr)\right)^{\!+}
    \label{eq:profile}
\end{equation}
where $(x)^+ = \max(0, x)$. The resulting profile
$\mathbf{r}_i = \bigl[r_{ik}^{(t)}\bigr]_{k \in K,\, t = 1
\ldots D_i}$ is a $D_i \times |K|$ matrix specifying the resource
consumption of activity~$i$ in each period of its execution.
\subsection{Scheduling Loop}
\label{subsec:ssgs_loop}
The actual scheduling works as a priority-driven loop. A residual
capacity matrix $\mathrm{rem}_{tk}$ tracks how much of each
resource is still available in each period, starting at the full
capacities~$R_k$.
At each step, the \textit{eligible set}~$\mathcal{E}$ collects all
unscheduled activities whose predecessors have already been placed:
\begin{equation}
    \mathcal{E} = \bigl\{\, i \in V \setminus \mathcal{S} :
    \forall\, j \in \mathrm{Pred}(i),\;
    j \in \mathcal{S} \,\bigr\}
    \label{eq:eligible}
\end{equation}
where $\mathcal{S}$ is the set of already-scheduled activities.
Among the eligible activities, the one with the smallest earliest
start time is picked, breaking ties by index:
\begin{equation}
    i^* = \arg\min_{i \in \mathcal{E}}\;
    \bigl(\mathrm{EST}_i,\; i\bigr)
    \label{eq:priority}
\end{equation}
Here $\mathrm{EST}_i$ is computed on the fly as the later of the
static earliest start~$ES_i$ from the precedence network and the
latest predecessor finish time:
$\mathrm{EST}_i = \max\!\bigl(ES_i,\;
\max_{j \in \mathrm{Pred}(i)} (S_j + D_j)\bigr)$.
The selected activity~$i^*$ is then placed at the first
time~$t_0 \geq \mathrm{EST}_{i^*}$ where enough residual capacity
exists in every period of its profile:
\begin{equation}
    S_{i^*} = \min\left\{\, t_0 \geq \mathrm{EST}_{i^*} :
    r_{i^* k}^{(\delta)} \leq \mathrm{rem}_{(t_0 + \delta),\, k},\;
    \forall\, \delta \in \{0, \ldots, D_{i^*}\!-\!1\},\;
    \forall\, k \in K \,\right\}
    \label{eq:ssgs_placement}
\end{equation}
After placement, the residual capacities are reduced:
$\mathrm{rem}_{(t_0 + \delta),\, k} \leftarrow
\mathrm{rem}_{(t_0 + \delta),\, k} - r_{i^* k}^{(\delta)}$ for
all $\delta$ and~$k$. The schedule
$S^{\mathrm{SSGS}} = (S_0, S_1, \ldots, S_{n+1})$ that comes out
of this loop is feasible by construction and is handed to the R\&F
phase as a warm-start. Algorithm~\ref{alg:ssgs} summarises the
complete procedure; Figure~\ref{fig:ssgs_flowchart} gives a visual
overview.
% ============================================================
% TikZ diagram: SSGS Flowchart for the SAA-FSRCPSP
% Layout adapted from Mendl & Rieck (Annals of OR), Fig. 1
% Content matches SSGS.py and Algorithm 1
%
% Include in Chapter 4 with:
%   % ============================================================
% TikZ diagram: SSGS Flowchart for the SAA-FSRCPSP
% Layout adapted from Mendl & Rieck (Annals of OR), Fig. 1
% Content matches SSGS.py and Algorithm 1
%
% Include in Chapter 4 with:
%   % ============================================================
% TikZ diagram: SSGS Flowchart for the SAA-FSRCPSP
% Layout adapted from Mendl & Rieck (Annals of OR), Fig. 1
% Content matches SSGS.py and Algorithm 1
%
% Include in Chapter 4 with:
%   \input{figures/ssgs_flowchart.tex}
%
% Required in preamble:
%   \usepackage{tikz}
%   \usetikzlibrary{shapes.geometric, arrows.meta, positioning,
%                   backgrounds, calc}
% ============================================================
\begin{figure}[ht]
\centering
\resizebox{\textwidth}{!}{%
\begin{tikzpicture}[
    node distance = 0.45cm and 0.55cm,
    %% Gray boxes – main scheduling steps
    gbox/.style   = {rectangle, draw=black!50, fill=gray!12,
                     minimum height=0.85cm, minimum width=2.2cm,
                     align=center, font=\footnotesize,
                     inner sep=4pt},
    %% Teal boxes – scenario / profile construction
    tbox/.style   = {rectangle, draw=teal!60, fill=teal!8,
                     minimum height=0.85cm, minimum width=2.2cm,
                     align=center, font=\footnotesize,
                     inner sep=4pt},
    %% Blue box – output
    bbox/.style   = {rectangle, draw=blue!50, fill=blue!8,
                     minimum height=0.85cm, minimum width=2.2cm,
                     align=center, font=\footnotesize,
                     inner sep=4pt},
    %% Intermediate value labels (between boxes)
    val/.style    = {font=\footnotesize, text=black!80},
    %% Start / end markers
    io/.style     = {font=\footnotesize},
    %% Arrows
    arr/.style    = {-{Stealth[length=4pt]}, draw=black!60},
    %% Small annotation labels
    lbl/.style    = {font=\scriptsize, text=black!55},
]

%% ── ROW 1: scheduling loop ────────────────────────────────
\node[io] (s0) {$S_0 = 0$};

\node[gbox, right=0.6cm of s0] (eligible)
    {Determine $\mathcal{E}$\\(eligible set)};

\node[gbox, right=of eligible] (select)
    {Select $j \in \mathcal{E}$\\(highest priority)};

\node[gbox, right=of select] (est)
    {Calculate time-\\feasible $\mathrm{EST}_j$};

\node[val, right=0.35cm of est] (tprime) {$t' = \mathrm{EST}_j$};

\node[tbox, right=0.35cm of tprime] (worst)
    {Select worst-case\\scenario $\hat{\pi}_j$};

\node[tbox, right=of worst] (profile)
    {Build resource\\profile $(D_j, \mathbf{r}_j)$};

\node[gbox, right=of profile] (resfeas)
    {Calculate resource-\\feasible $t^* \geq \mathrm{EST}_j$};

%% ── Loop-back arrow (above, like Annals Fig. 1) ──────────
\draw[arr]
    (resfeas.north) -- ++(0, 0.8)
    -| (eligible.north)
    node[lbl, pos=0.15, above] {$\bar{C} \neq \emptyset$};

%% ── ROW 1 arrows ─────────────────────────────────────────
\draw[arr] (s0)       -- (eligible);
\draw[arr] (eligible) -- (select);
\draw[arr] (select)   -- (est);
\draw[arr] (est)      -- (tprime);
\draw[arr] (tprime)   -- (worst);
\draw[arr] (worst)    -- (profile);
\draw[arr] (profile)  -- (resfeas);

%% ── Exit: C̄ = ∅ → output ────────────────────────────────
\node[val, below=1.0cm of resfeas] (sval)
    {$S^{\mathrm{SSGS}}$};
\node[bbox, left=0.5cm of sval] (output)
    {Return\\$x_{\mathrm{init}} = S^{\mathrm{SSGS}}$};
\node[io, left=0.6cm of output] (xinit) {$x_{\mathrm{init}}$};

\draw[arr] (resfeas.south) -- (sval.north)
    node[lbl, pos=0.5, right] {$\bar{C} = \emptyset$};
\draw[arr] (sval) -- (output);
\draw[arr] (output) -- (xinit);

%% ── Gray dashed background for scheduling loop ───────────
\begin{scope}[on background layer]
    \node[draw=black!25, dashed, rounded corners=4pt,
          inner xsep=8pt, inner ysep=16pt, fill=gray!4,
          fit=(s0)(eligible)(select)(est)(tprime)
              (worst)(profile)(resfeas)]
          (graybox) {};
\end{scope}

\end{tikzpicture}
}
\caption{Serial schedule generation scheme for the SAA-FSRCPSP.}
\label{fig:ssgs_flowchart}
\end{figure}

%
% Required in preamble:
%   \usepackage{tikz}
%   \usetikzlibrary{shapes.geometric, arrows.meta, positioning,
%                   backgrounds, calc}
% ============================================================
\begin{figure}[ht]
\centering
\resizebox{\textwidth}{!}{%
\begin{tikzpicture}[
    node distance = 0.45cm and 0.55cm,
    %% Gray boxes – main scheduling steps
    gbox/.style   = {rectangle, draw=black!50, fill=gray!12,
                     minimum height=0.85cm, minimum width=2.2cm,
                     align=center, font=\footnotesize,
                     inner sep=4pt},
    %% Teal boxes – scenario / profile construction
    tbox/.style   = {rectangle, draw=teal!60, fill=teal!8,
                     minimum height=0.85cm, minimum width=2.2cm,
                     align=center, font=\footnotesize,
                     inner sep=4pt},
    %% Blue box – output
    bbox/.style   = {rectangle, draw=blue!50, fill=blue!8,
                     minimum height=0.85cm, minimum width=2.2cm,
                     align=center, font=\footnotesize,
                     inner sep=4pt},
    %% Intermediate value labels (between boxes)
    val/.style    = {font=\footnotesize, text=black!80},
    %% Start / end markers
    io/.style     = {font=\footnotesize},
    %% Arrows
    arr/.style    = {-{Stealth[length=4pt]}, draw=black!60},
    %% Small annotation labels
    lbl/.style    = {font=\scriptsize, text=black!55},
]

%% ── ROW 1: scheduling loop ────────────────────────────────
\node[io] (s0) {$S_0 = 0$};

\node[gbox, right=0.6cm of s0] (eligible)
    {Determine $\mathcal{E}$\\(eligible set)};

\node[gbox, right=of eligible] (select)
    {Select $j \in \mathcal{E}$\\(highest priority)};

\node[gbox, right=of select] (est)
    {Calculate time-\\feasible $\mathrm{EST}_j$};

\node[val, right=0.35cm of est] (tprime) {$t' = \mathrm{EST}_j$};

\node[tbox, right=0.35cm of tprime] (worst)
    {Select worst-case\\scenario $\hat{\pi}_j$};

\node[tbox, right=of worst] (profile)
    {Build resource\\profile $(D_j, \mathbf{r}_j)$};

\node[gbox, right=of profile] (resfeas)
    {Calculate resource-\\feasible $t^* \geq \mathrm{EST}_j$};

%% ── Loop-back arrow (above, like Annals Fig. 1) ──────────
\draw[arr]
    (resfeas.north) -- ++(0, 0.8)
    -| (eligible.north)
    node[lbl, pos=0.15, above] {$\bar{C} \neq \emptyset$};

%% ── ROW 1 arrows ─────────────────────────────────────────
\draw[arr] (s0)       -- (eligible);
\draw[arr] (eligible) -- (select);
\draw[arr] (select)   -- (est);
\draw[arr] (est)      -- (tprime);
\draw[arr] (tprime)   -- (worst);
\draw[arr] (worst)    -- (profile);
\draw[arr] (profile)  -- (resfeas);

%% ── Exit: C̄ = ∅ → output ────────────────────────────────
\node[val, below=1.0cm of resfeas] (sval)
    {$S^{\mathrm{SSGS}}$};
\node[bbox, left=0.5cm of sval] (output)
    {Return\\$x_{\mathrm{init}} = S^{\mathrm{SSGS}}$};
\node[io, left=0.6cm of output] (xinit) {$x_{\mathrm{init}}$};

\draw[arr] (resfeas.south) -- (sval.north)
    node[lbl, pos=0.5, right] {$\bar{C} = \emptyset$};
\draw[arr] (sval) -- (output);
\draw[arr] (output) -- (xinit);

%% ── Gray dashed background for scheduling loop ───────────
\begin{scope}[on background layer]
    \node[draw=black!25, dashed, rounded corners=4pt,
          inner xsep=8pt, inner ysep=16pt, fill=gray!4,
          fit=(s0)(eligible)(select)(est)(tprime)
              (worst)(profile)(resfeas)]
          (graybox) {};
\end{scope}

\end{tikzpicture}
}
\caption{Serial schedule generation scheme for the SAA-FSRCPSP.}
\label{fig:ssgs_flowchart}
\end{figure}

%
% Required in preamble:
%   \usepackage{tikz}
%   \usetikzlibrary{shapes.geometric, arrows.meta, positioning,
%                   backgrounds, calc}
% ============================================================
\begin{figure}[ht]
\centering
\resizebox{\textwidth}{!}{%
\begin{tikzpicture}[
    node distance = 0.45cm and 0.55cm,
    %% Gray boxes – main scheduling steps
    gbox/.style   = {rectangle, draw=black!50, fill=gray!12,
                     minimum height=0.85cm, minimum width=2.2cm,
                     align=center, font=\footnotesize,
                     inner sep=4pt},
    %% Teal boxes – scenario / profile construction
    tbox/.style   = {rectangle, draw=teal!60, fill=teal!8,
                     minimum height=0.85cm, minimum width=2.2cm,
                     align=center, font=\footnotesize,
                     inner sep=4pt},
    %% Blue box – output
    bbox/.style   = {rectangle, draw=blue!50, fill=blue!8,
                     minimum height=0.85cm, minimum width=2.2cm,
                     align=center, font=\footnotesize,
                     inner sep=4pt},
    %% Intermediate value labels (between boxes)
    val/.style    = {font=\footnotesize, text=black!80},
    %% Start / end markers
    io/.style     = {font=\footnotesize},
    %% Arrows
    arr/.style    = {-{Stealth[length=4pt]}, draw=black!60},
    %% Small annotation labels
    lbl/.style    = {font=\scriptsize, text=black!55},
]

%% ── ROW 1: scheduling loop ────────────────────────────────
\node[io] (s0) {$S_0 = 0$};

\node[gbox, right=0.6cm of s0] (eligible)
    {Determine $\mathcal{E}$\\(eligible set)};

\node[gbox, right=of eligible] (select)
    {Select $j \in \mathcal{E}$\\(highest priority)};

\node[gbox, right=of select] (est)
    {Calculate time-\\feasible $\mathrm{EST}_j$};

\node[val, right=0.35cm of est] (tprime) {$t' = \mathrm{EST}_j$};

\node[tbox, right=0.35cm of tprime] (worst)
    {Select worst-case\\scenario $\hat{\pi}_j$};

\node[tbox, right=of worst] (profile)
    {Build resource\\profile $(D_j, \mathbf{r}_j)$};

\node[gbox, right=of profile] (resfeas)
    {Calculate resource-\\feasible $t^* \geq \mathrm{EST}_j$};

%% ── Loop-back arrow (above, like Annals Fig. 1) ──────────
\draw[arr]
    (resfeas.north) -- ++(0, 0.8)
    -| (eligible.north)
    node[lbl, pos=0.15, above] {$\bar{C} \neq \emptyset$};

%% ── ROW 1 arrows ─────────────────────────────────────────
\draw[arr] (s0)       -- (eligible);
\draw[arr] (eligible) -- (select);
\draw[arr] (select)   -- (est);
\draw[arr] (est)      -- (tprime);
\draw[arr] (tprime)   -- (worst);
\draw[arr] (worst)    -- (profile);
\draw[arr] (profile)  -- (resfeas);

%% ── Exit: C̄ = ∅ → output ────────────────────────────────
\node[val, below=1.0cm of resfeas] (sval)
    {$S^{\mathrm{SSGS}}$};
\node[bbox, left=0.5cm of sval] (output)
    {Return\\$x_{\mathrm{init}} = S^{\mathrm{SSGS}}$};
\node[io, left=0.6cm of output] (xinit) {$x_{\mathrm{init}}$};

\draw[arr] (resfeas.south) -- (sval.north)
    node[lbl, pos=0.5, right] {$\bar{C} = \emptyset$};
\draw[arr] (sval) -- (output);
\draw[arr] (output) -- (xinit);

%% ── Gray dashed background for scheduling loop ───────────
\begin{scope}[on background layer]
    \node[draw=black!25, dashed, rounded corners=4pt,
          inner xsep=8pt, inner ysep=16pt, fill=gray!4,
          fit=(s0)(eligible)(select)(est)(tprime)
              (worst)(profile)(resfeas)]
          (graybox) {};
\end{scope}

\end{tikzpicture}
}
\caption{Serial schedule generation scheme for the SAA-FSRCPSP.}
\label{fig:ssgs_flowchart}
\end{figure}

\begin{algorithm}[ht]
\caption{Procedure of the proposed serial schedule generation scheme}
\label{alg:ssgs}
\begin{algorithmic}[1]
\State Initialise $S_0 \leftarrow 0$, $D_0 \leftarrow 0$,
       $\mathcal{S} \leftarrow \{0\}$, and
       $\mathrm{rem}_{tk} \leftarrow R_k$ $\forall\, t \in T$,
       $k \in K$
\While{$\mathcal{S} \neq V$}
    \State Determine eligible set $\mathcal{E}$ containing all
           unscheduled activities with
           $\mathrm{Pred}(j) \subseteq \mathcal{S}$
    \State Select $j \in \mathcal{E}$ with highest priority
           \Comment{e.g., EST rule}
    \State Compute $\mathrm{EST}_j \leftarrow
           \max\bigl(ES_j,\;
           \max_{i \in \mathrm{Pred}(j)} (S_i + D_i)\bigr)$
    \State Select worst-case scenario
           $\hat{\pi} \leftarrow \arg\max_{\pi \in \Pi}\;
           \max_{k \in K}\;
           \lceil w_{jk}^{\pi} / \overline{r}_{jk} \rceil$
           \Comment{Step~1}
    \State Set $\hat{w}_k \leftarrow w_{jk}^{\hat{\pi}}$
           $\forall\, k \in K$
    \State Compute $D_j \leftarrow \max_{k \in K}\;
           \lceil \hat{w}_k / \overline{r}_{jk} \rceil$
           \Comment{Step~2}
    \For{each $k \in K$}
           \Comment{front-loaded profile}
        \State Assign $\underline{r}_{jk}$ $\forall\,
               \delta = 0, \ldots, D_j\!-\!1$; distribute
               remaining $\hat{w}_k - D_j \cdot
               \underline{r}_{jk}$ up to $\overline{r}_{jk}$
    \EndFor
    \State Find earliest $t^* \geq \mathrm{EST}_j$ such that
           $r_{jk}^{(\delta)} \leq
           \mathrm{rem}_{(t^*\!+\delta),k}$
           $\forall\, \delta, k$
           \Comment{Step~3}
    \State $S_j \leftarrow t^*$; update
           $\mathrm{rem}_{(t^*\!+\delta),k} \leftarrow
           \mathrm{rem}_{(t^*\!+\delta),k} -
           r_{jk}^{(\delta)}$ $\forall\, \delta, k$
    \State $\mathcal{S} \leftarrow \mathcal{S} \cup \{j\}$
\EndWhile
\State \Return $S^{\mathrm{SSGS}}$
\end{algorithmic}
\end{algorithm}
% ------------------------------------------------------------
\section{The Relax-and-Fix Decomposition}
\label{sec:baseline_rf}
% ------------------------------------------------------------
With the SSGS solution~$S$ in hand, the R\&F phase tries to
improve each activity's start time by solving a sequence of MIP
subproblems. Activities are visited in topological order
$i = 0, 1, \ldots, |V|-1$. Two schedules are maintained: a
\textit{working schedule}~$S'$, updated whenever a subproblem
yields a feasible solution, and a \textit{best schedule}~$S^*$,
which keeps the lowest makespan found so far. Both start as copies
of the SSGS solution.
\subsection{Time Window Construction}
\label{subsec:time_window}
Each activity~$i$ gets a restricted time window around its current
start time~$S'_i$:
\begin{equation}
    \widetilde{W}_i =
    \left[\,\max\{ES_i,\; S'_i - w\},\;
           \min\{LS_i,\; S'_i + w\}\,\right]
    \label{eq:time_window}
\end{equation}
where $w$ is drawn uniformly from $\{1, \ldots, w_{\max}\}$
($w_{\max} = 10$ in this thesis). Only start times inside
$\widetilde{W}_i$ are considered; all variables $x_{it}$ with
$t \notin \widetilde{W}_i$ are fixed to zero. Drawing~$w$ randomly
keeps the algorithm from getting stuck in one neighbourhood size
and lets different passes explore different regions of the solution
space \parencite{pochet2006}.
\subsection{Variable States}
\label{subsec:variable_states}
Each subproblem assigns the decision variables $x_{it}$ one of
three states (see also Figure~\ref{fig:baseline_pipeline}):
\begin{itemize}
    \item \textbf{Fixed:} Activities $j < i$ that have already
          been processed. Their variables are locked:
          $x_{j, S'_j} = 1$, $x_{jt} = 0$ for $t \neq S'_j$.
    \item \textbf{Binary:} The current activity~$i$, kept as an
          integer variable within its time window:
          $x^B_{it} \in \{0,1\}$ for $t \in \widetilde{W}_i$.
    \item \textbf{Relaxed:} All later activities $j > i$, whose
          binary variables are relaxed to continuous:
          $x^C_{jt} \in [0,1]$ for all $t \in T$.
\end{itemize}
Because only one activity's variables are binary at a time, the
subproblem is much easier to solve than the full MIP, yet the
precedence and resource structure of the SAA-FSRCPSP is still
present through the continuous relaxation of future activities.
\subsection{Subproblem Solution and Acceptance}
\label{subsec:subproblem}
Each subproblem is the SAA-FSRCPSP with the variable states above,
solved by Gurobi \parencite{gurobi2024} within a per-model time
limit~$T_{\mathrm{model}}$. Three outcomes are possible:
\begin{enumerate}
    \item \textbf{Feasible solution found.} A candidate
          schedule~$S'_{\mathrm{cand}}$ is extracted. Activity~$i$
          is fixed at~$S'_{\mathrm{cand},i}$ and the working
          schedule is updated ($S' \leftarrow S'_{\mathrm{cand}}$),
          regardless of whether the makespan improved. If the
          makespan did improve
          ($C_{\max}(S'_{\mathrm{cand}}) < C_{\max}(S^*)$), the
          best schedule is updated too.
    \item \textbf{Infeasible, retries remain.} Sometimes the time
          window is too tight and the subproblem has no feasible
          solution. In that case a new~$w$ is drawn and the
          subproblem is re-solved. Up to $z_{\max} = 2$ attempts are
          made per activity.
    \item \textbf{Infeasible, retries exhausted.} If all
          $z_{\max}$ attempts fail, activity~$i$ falls back to its
          SSGS start time:
          $S'_i \leftarrow S^{\mathrm{SSGS}}_i$. This way the
          algorithm can always continue, since the SSGS solution is
          feasible by construction.
\end{enumerate}
One more check happens after every subproblem: if the solution
turns out to be \textit{integral}, meaning all effective start time
variables $x^{eff}_{it}$ are binary across all activities, the
algorithm stops immediately. An integral solution means the current
schedule is feasible for the full SAA-FSRCPSP with no LP relaxation
gaps, so there is no point in continuing.
\subsection{Stopping Criteria}
\label{subsec:stopping}
The outer loop repeats the full pass over all activities until one
of two things happens:
\begin{enumerate}
    \item The wall-clock time limit~$T_{\max}$ is reached.
    \item An integral solution is found.
    % \item \texttt{max\_no\_improve} consecutive passes go by
    %       without improving~$C_{\max}(S^*)$. (The default is~0,
    %       which disables this criterion.)
\end{enumerate}
Algorithm~\ref{alg:baseline} summarises the complete R\&F
procedure.
\begin{algorithm}[ht]
\caption{Procedure of the proposed activity-oriented F\&R}
\label{alg:baseline}
\begin{algorithmic}[1]
\State \mbox{Determine a feasible greedy solution $S$ by applying a serial schedule generation scheme}
%\State Check feasibility of $S$ over all $\pi \in \mathcal{S}$ with regard to $(1-\epsilon)$
\While{stop criterion is not met}
    \For{$i=0$ \textbf{to} $|V|-1$}
        \State Choose parameter $w = \mathrm{rand}(1,10)$
        % \State Choose parameter $w_i = \max\!\left(2,\left\lfloor 0.15\,(LS_i-ES_i)\right\rfloor\right)$
        \State Define restricted time window $\widetilde{W}_i = [\max\{ES_i,S_i-w_i\},\,\min\{LS_i,S_i+w_i\}]$
        \State Fix $x_{it}= 0$ $\forall$ $t \notin \widetilde{W}_i$
        \State \parbox[t]{\dimexpr\linewidth-\algorithmicindent\relax}{%
        Solve subproblem of the SAA-FSRCPSP using $x_{it}\in\{0,1\}$, $\forall$ $t\in\widetilde{W}_i$, $x_{jt}\in[0,1]$, $\forall$ $j>i,\ t\in T$, and all already fixed decision variables%
}
        \If{solution $S'$ of the subproblem is feasible}
            \State Fix $x_{iS'_i} = 1$ and $x_{it} = 0$ $, \forall$ $t\in\widetilde{W}_i\setminus\{S'_i\}$
            \If{$S'$ improves the makespan}
                \State Update $S = S'$ with $S_{n+1} = S'_{n+1}$
            \EndIf
        \Else
            \State $i = i-1$ and choose another $w$
        \EndIf
    \EndFor
\EndWhile
\State \Return $S$
\end{algorithmic}
\end{algorithm}
% ------------------------------------------------------------
\section{Complexity and Limitations}
\label{sec:baseline_limits}
% ------------------------------------------------------------
Each pass through the activity list solves up to $|V| = n + 2$
subproblems. A single subproblem has one binary activity and
$O(n \cdot T \cdot |S|)$ continuous variables for the relaxed
ones, so its size grows linearly with the number of remaining
activities. To give a sense of scale: for $n = 30$, $|K| = 4$,
$|S| = 10$, the full model has roughly 250{,}000--319{,}000
variables and 499{,}000--638{,}000 constraints
(cf.\ Table~\ref{tab:model_size}); for $n = 60$ these numbers
exceed one million.
Even with the per-subproblem limit~$T_{\mathrm{model}}$, the
wall-clock time adds up. At $n = 30$, $|K| = 4$, one pass means up
to 32~subproblems $\times$ 60~seconds each, so about 32~minutes in
the worst case, and the outer loop may run several passes.
The main structural limitation of this baseline is that it fixes
activities \textit{one at a time} in index order: the solver sees
only a single binary variable per subproblem, producing weak LP
relaxations, and the processing order carries no information
about which activities are related \parencite{helber2010}.
Chapter~\ref{chap:grouprf} addresses both issues.
